\cgeentry{cge_optimize}{Choose policy instruments to maximise an outcome or reach targets}

\cgesec{Description}

\cmd{cge_optimize} chooses the values of a few policy instruments to maximise
or minimise a report measure, or to reach targets for report measures,
solving the model at every trial value. An instrument is one element of a
parameter that \manref{cge_shock} changes, such as a tariff rate, or one
element of a variable that the closure fixes, such as a tax adjustment
factor. The optimisation starts from the instruments' current values, with
any shocks already applied, under the current closure. The optimum is stored
as a simulation, and the instruments are left at their optimal values as
current shocks, so that \mbox{\manref{cge_report}} and \mbox{\manref{cge_results}} show it.

\cgesec{Quick start}

The tariff on manufactures, between $-0.5$ and 1, that maximises the
equivalent variation of household \cmd{hhd}
\begin{stcode}
cge_optimize, instruments(tm(com_man)) maximize(welfare equivalent_variation hhd) ///
    lower(tm(com_man):-0.5) upper(tm(com_man):1)
\end{stcode}
Both tariffs that maximise it while tariff revenue stays at 8, storing the
optimum as simulation \cmd{two}
\begin{stcode}
cge_optimize, instruments(tm(com_agr) tm(com_man)) maximize(welfare equivalent_variation hhd) ///
    targets(government tax_import = 8) name(two)
\end{stcode}
The direct tax adjustment factor that keeps government savings at 3, with the
iterations displayed
\begin{stcode}
cge_optimize, instruments(TINSADJ) targets(government savings = 3) trace
\end{stcode}
Government purchases of health services and the direct tax factor that reach
targets for real GDP and government savings
\begin{stcode}
cge_optimize, instruments(qg(com_hlt) TINSADJ) targets(macro gdp_real = 310, government savings = 3)
\end{stcode}

\cgesec{Syntax}

\begin{stsyntax}
\cmdb{cge_optimize}\texttt{,} \optarg{inst}{ruments}{list}
\{\optarg{max}{imize}{measure} $|$ \optarg{min}{imize}{measure} $|$
\optarg{targ}{ets}{spec}\} [\argu{options}]
\end{stsyntax}

\begin{opttable}
\optgroup{Main}
\optarg{inst}{ruments}{list} & the instruments; required\\
\optarg{max}{imize}{measure} & report measure to maximise\\
\optarg{min}{imize}{measure} & report measure to minimise\\
\optarg{targ}{ets}{spec} & targets for report measures\\
\optarg{low}{er}{spec} & lower bounds of instruments\\
\optarg{up}{per}{spec} & upper bounds of instruments\\
\optarg{}{name}{name} & name of the simulation that holds the optimum; default \cmd{optimum}\\
\optgroup{Solver}
\optarg{tol}{erance}{\#} & tolerance of each solve of the model; default \texttt{1e-12}\\
\optarg{iter}{ate}{\#} & maximum number of Newton iterations of each solve; default \texttt{100}\\
\opt{}{trace} & display the iterations of the optimiser\\
\end{opttable}

A \argu{measure} is a measure of \manref{cge_report}, written
\argu{table measure} [\argu{element}], for example
\cmd{welfare equivalent_variation hhd} or \cmd{macro gdp_real}. An objective,
targets or both must be given; \cmd{maximize()} and \cmd{minimize()} may not
be combined.

\cgesec{Options}

\begin{optdesc}
\item[\optarg{inst}{ruments}{list}] lists the instruments, separated by
spaces. An instrument is
\begin{itemize}\setlength{\itemsep}{0pt}\setlength{\labelsep}{0.6em}
\item one element of a parameter that \manref{cge_shock} changes, written
\argu{parameter}\cmd{(}\argu{element}\cmd{)}, such as \cmd{tm(com_man)} or
\cmd{qg(com_hlt)}; the element may be left out when the parameter has only
one;
\item a scalar parameter, such as \cmd{trgr};
\item one element of a user parameter (\manref{cge_parameter}), such as
\cmd{w(com_agr)} or \cmd{sh(com_hlt,hhd)}, or a scalar one; or
\item one element of a variable that the closure fixes, such as \cmd{GADJ},
\cmd{TINSADJ} or \cmd{QFS(lab)}. A variable that the closure does not fix is
an error; it can be fixed by a swap (\manref{cge_closure}, \cmd{swap()}).
\end{itemize}
An instrument may be given only once.

\item[\optarg{max}{imize}{measure}] gives the measure to maximise.

\item[\optarg{min}{imize}{measure}] gives the measure to minimise.

\item[\optarg{targ}{ets}{spec}] gives targets, separated by commas, as
\argu{table measure} [\argu{element}] \cmd{=} \argu{value}, for example
\cmd{targets(government savings = 3, macro investment_real = 50)}. Without an
objective there must be one target for each instrument. With an objective,
the targets are equality constraints.

\item[\optarg{low}{er}{spec}] and \optarg{up}{per}{spec} give bounds, as
\argu{instrument}\cmd{:}\argu{value} pairs with the instrument written as in
\cmd{instruments()}, for example
\cmd{lower(tm(com_man):0) upper(tm(com_man):0.2)}. An instrument without a
bound is unbounded on that side. A lower bound must be below the upper
bound, and the current value of each instrument must lie strictly within its
bounds.

\item[\optarg{}{name}{name}] names the simulation that holds the optimum;
the default is \cmd{optimum}. A simulation of that name is replaced.

\item[\optarg{tol}{erance}{\#}] sets the convergence tolerance of each solve
of the model; the default is \texttt{1e-12}.

\item[\optarg{iter}{ate}{\#}] sets the maximum number of Newton iterations of
each solve of the model; the default is 100.

\item[\opt{}{trace}] displays the objective and the largest gradient of the
free instruments at each iteration, or, with targets alone, the largest
scaled gap from the targets.
\end{optdesc}

\cgesec{Remarks and examples}

\cmd{cge_optimize} is meant for a few instruments: a tax rate, the volume of
a public service, a transfer. At each trial value of the instruments the
model is solved by Newton's method, from the solution of the trial before,
and the measures are computed at the solution. The objective and the targets
are values of the measures, such as amounts in the units of the SAM, not
percentage changes from the base; the equivalent and compensating variations
are amounts of base-price spending, not per cent of base consumption.
\Chapref{chap:optimize} states the problem and the algorithm.

\cgesubsec{Example 1: a welfare-maximising tariff}

The synthetic SAM is calibrated as in \manref{cge_solve}. The base values of
government revenue and its components, and a closure under which government
savings are fixed and the direct tax rates adjust, are
\stlog{opt_base}
The tariff on manufactures that maximises the equivalent variation of the
household, within $[-0.5, 1]$, with the iterations displayed:
\stlog{opt_max}
The objective is the equivalent variation in the units of the SAM; it is zero
at the start, the base. The gradient falls to $7.11 \times 10^{-9}$ in five
Newton iterations, and the 32 solves of the model include those of the
finite differences. On the synthetic SAM the optimal tariff is about 30 per
cent, above the base rate of 10 per cent. The result depends on the SAM and
the elasticities, which here are synthetic and illustrative, and is not a
finding about any economy.

The optimum is stored as simulation \cmd{optimum} and left as the current
shock:
\stlog{opt_return}
The equivalent variation of 0.768 is 0.334 per cent of base consumption, as
\manref{cge_report} shows it. Government savings are fixed by the closure:
tariff revenue more than doubles and direct taxes fall by 38 per cent.

\cgesubsec{Example 2: an objective with a target}

Both tariffs are chosen to maximise welfare while tariff revenue stays at its
base value of 8. The target is an equality constraint, met by the augmented
Lagrangian method:
\stlog{opt_revenue}
The optimum subsidises imports of agriculture slightly and taxes those of
manufactures at 16 per cent. Each round of the augmented Lagrangian method
maximises again, so the optimisation took 419 solves of the model.

\cgesubsec{Example 3: targets alone}

With as many targets as instruments and no objective, the instruments are
found by Newton's method on the targets. After the removal of every tariff,
under a closure in which government savings adjust, the direct tax
adjustment factor that keeps government savings at their base value of 3,
found in two Newton iterations, is
\stlog{opt_targets}
Holding government savings fixed with the direct tax rates adjusting is what
the closure \cmd{fixgsav} does. Solving the same shock under that closure
gives the same factor:
\stlog{opt_fixgsav}

\cgesubsec{Bounds, failures and local optima}

An instrument at a bound whose gradient points outwards is held there, and
the display marks it \cmd{(at its bound)}. A trial value at which the model
does not solve counts in \res{r(N_failed)}; the next trial starts from the
last solution, and the optimiser steps back. The model must solve at the
starting values. Derivatives are found by central differences of the solved
model, and a smooth objective is flat at its optimum, so the optimum is
located to about $10^{-6}$ relative to the instruments' scale. An optimum may
be local: other starting values, set with \manref{cge_shock}, can be tried.

\cgesec{Stored results}

\cmd{cge_optimize} stores the following in \cmd{r()}:
\begin{storedres}
\resgroup{Scalars}
\res{r(converged)} & 1 if the optimisation converged, 0 otherwise\\
\res{r(objective)} & the objective at the optimum, with an objective\\
\res{r(N_evaluations)} & number of solves of the model\\
\res{r(N_failed)} & number of trial values at which the model did not solve\\
\res{r(iterations)} & number of iterations of the optimiser\\
\resgroup{Macros}
\res{r(name)} & name of the simulation that holds the optimum\\
\res{r(instruments_list)} & the instruments\\
\resgroup{Matrices}
\res{r(instruments)} & the instruments at the optimum\\
\res{r(start)} & their starting values\\
\res{r(targets)} & the target measures at the optimum, with targets\\
\end{storedres}

\cgesec{Methods and formulas}

With instruments $u$, the model's solution $x(u)$ under the current closure,
an objective measure $f$ and target measures $g_j$,
\begin{equation*}
\max_u\; f\bigl(x(u), u\bigr) \quad\text{subject to}\quad
g_j\bigl(x(u), u\bigr) = \bar g_j, \qquad \underline u \le u \le \overline u,
\end{equation*}
with $-f$ for \cmd{minimize()}. The instruments are scaled as $u = u^0 + s
\odot z$, with $u^0$ the starting values and $s_k = |u^0_k|$, or 1 when
$|u^0_k| \le 10^{-12}$. The target gaps are scaled as
$r_j = (g_j - \bar g_j)/\tau_j$, with $\tau_j = |\bar g_j|$, or 1 when
$|\bar g_j| \le 10^{-12}$.

\textit{Targets alone.} Newton's method solves $r(z) = 0$. The Jacobian comes
from central differences with steps $h_k = 10^{-5}\max(1, |z_k|)$; the step
$\Delta z = -J^{-1} r$ is projected on the bounds and halved, up to 30 times,
until $\max_j |r_j|$ falls. The iterations stop, converged, when
$\max_j |r_j| < 10^{-11}$; they stop without converging when a trial
value does not solve, when 30 halvings do not reduce the gap, and after 50
iterations.

\textit{An objective.} A projected Newton method maximises
\begin{equation*}
F(z) = \frac{f}{\max(1, |f^0|)} - \lambda' r - \frac{\rho}{2}\, r' r,
\end{equation*}
with $f^0$ the objective at the start; without targets, $F$ is the scaled
objective alone. The gradient comes from central differences with step
$10^{-5}$ (one-sided at a bound); an instrument at a bound whose gradient
points outwards is held there. The iterations converge when the largest
gradient of the free instruments is below $10^{-8}$. The Hessian of the free
instruments comes from central second differences with step $10^{-3}$. The
step is $-H^{-1}\nabla F$ when $H$ is negative definite, and the gradient,
divided by $\max(1, \max_k |\partial F/\partial z_k|)$, otherwise; it is
projected on the bounds and halved, up to 40 times, until $F$ rises. If no
halving raises $F$, the iterations stop, and they are counted as converged
if the largest free gradient is below $10^{-6}$; they also stop, converged,
when a step changes $z$ by less than $10^{-12}$, and after at most 100
iterations.

\textit{Targets with an objective.} By the augmented Lagrangian method, with
$\lambda = 0$ and $\rho = 100$ at the start, each round maximises $F$ and
then sets $\lambda \leftarrow \lambda + \rho\,r$; $\rho$ is multiplied by 10
when $\max_j |r_j|$ has fallen by less than three quarters since the round
before. The rounds stop, converged, when $\max_j |r_j| < 10^{-10}$, and after
at most 40 rounds.

\Chapref{chap:optimize} gives the details.

\cgesec{Also see}

\begin{alsosee}
\seeentry{cge_shock}{Change a parameter or a fixed variable}
\seeentry{cge_closure}{Choose the model closure}
\seeentry{cge_solve}{Solve the model}
\seeentry{cge_report}{Standard report tables for one or more simulations}
\seechapter{chap:optimize}{Policy optimisation}
\end{alsosee}
